% AAAI 2027 main-track submission — skeleton.
% Source-of-truth narrative: paper/visualization-tools-putnam.md
% Operating plan:             paper/publication_plan.md
% Pre-registration:           paper/pre_registration.md
%
% NOTE: the AAAI 2027 author kit (`aaai27.sty`, `aaai27.bst`) is not
% yet available at the time of authoring (2026-05-21). When released,
% drop the files into `aaai-style/` and switch the documentclass /
% packages below. Until then the article-class fallback lets us draft.

\documentclass[letterpaper]{article}
% \usepackage{aaai27}   % uncomment when style file is in place
\usepackage{times}
\usepackage{helvet}
\usepackage{courier}
\usepackage[hyphens]{url}
\usepackage{graphicx}
\usepackage{booktabs}
\usepackage{amsmath, amssymb}
\usepackage{microtype}

\frenchspacing
\setlength{\pdfpagewidth}{8.5in}
\setlength{\pdfpageheight}{11in}

\title{Discovery dominates value: \\
  why advertised visualization tools go unused by an autonomous LLM
  agent, and what to do about it}

% Author block — replace with the AAAI author block when submitting.
\author{
  Silvere Gangloff
}

\begin{document}
\maketitle

% ----------------------------------------------------------------------
\begin{abstract}
% Target: 200 words. Lead with the decomposition; instantiate it.
%
% Draft (to be tightened):
We study whether giving an autonomous LLM coding agent access to
visualization tools changes what it can solve. On a corpus of
statement-level Lean~4 formalizations of Putnam real-analysis
problems and a sandboxed harness, an A/B sweep finds that an
unattended \texttt{--print}-mode agent invokes the visualization tool
\emph{exactly zero times} across 147 runs, including on the
integral / asymptotic subset most amenable to a picture. We attribute
this to a general
\emph{discovery--value decomposition}: an advertised but optional
agent tool affects end-task performance by
$\text{tool\_value} \times P(\text{adopted} \mid \text{prompt, mode})$,
and in unattended mode the adoption term can be identically zero.
A forced-loop ablation lifts adoption to $100\%$ without changing
solve rates; a zero-tool prose-CoT control matches forced-viz on the
same corpus, showing the test set has no headroom for any tool. An
adversarial false-claim study finds that required visualization,
required SymPy evaluation, and required tool-free prose reflection
catch deliberately mis-stated answers at indistinguishable rates: the
corrective work is the required verification step, not its modality.
A constructive remedy --- returning inline content instead of file
paths --- lifts voluntary adoption by $X\%$. Implications for
agent-tool evaluation design are discussed.
\end{abstract}

% ----------------------------------------------------------------------
\section{Introduction}
% TODO: 1 page.
% - Hook (1 par): humans plot first; does an LLM agent?
% - Negative result (1 par): the answer is "no in autonomous mode."
% - The decomposition (1 par): tool_value x P(adopted | prompt, mode).
%   Name it. State it. This is the framing contribution.
% - Contributions (4 bullets): decomposition; empirical landscape;
%   constructive recommendation (inline content); methodological
%   point about CoT-incompetence baselines.

\section{The discovery--value decomposition}
% Crisp statement of the formula + assumptions. ~0.5 page.
% Position as a methodological contribution that survives outside the
% paper's specific empirics.

\section{Setup}
% Corpus, model, harness, judge, scoring. ~0.75 page.
% Refer to appendix for full details on Studies 1--5.

\section{Studies 6 \& 7 \& 8: power, design, modality}
% This is the empirical spine. ~3 pages.
%
% 6.1 Scaled adversarial (Study 6): modality equivalence on analytic
%     corpus at n=24/condition. Result table.
% 6.2 Inline-content design ablation (Study 7): adoption lift, with
%     hand-coded grounding analysis. Headline figure.
% 6.3 Geometric adversarial (Study 8): does picture beat numerics on
%     a corpus designed to favor pictures? Headline figure.

\section{Discussion}
% ~0.75 page.
% - What the decomposition implies for benchmarks.
% - Why "modality doesn't matter when verification is required" is
%   the right reading of Studies 6 and 8.
% - The MCP-tool-design lesson from Study 7.

\section{Threats to validity}
% ~0.25 page.
% - Single model (Claude family). Be explicit; cite the existence
%   case argument from the pre-reg.
% - LLM judge: cite kappa from human audit in appendix.
% - N, corpus narrowness.

\section{Related work}
% ~0.75 page.
% - ReAct, Toolformer, ToolLLM (tool-use evaluation).
% - MM-Math, MathVista (mathematical visual reasoning).
% - GAIA, AgentBench (autonomous-agent benchmarks).
% - Tool-use eval workshops 2024--2025.
% - Position contribution: most prior work measures end-task
%   performance with advertised tools; we expose the adoption
%   confound.

\section{Conclusion}
% ~0.25 page.
% One paragraph restating decomposition + one-line summary of each
% empirical contribution + the constructive recommendation.

% Bibliography
% \bibliographystyle{aaai27}     % uncomment with style file
\bibliographystyle{plain}
\bibliography{references}

\end{document}
